\section{CNN 经典架构}

\subsection{ImageNet 竞赛与架构演进}

\begin{figure}[H]
\centering
\includegraphics[width=0.95\textwidth]{slides-images/slide-020.jpg}
\caption{ImageNet 竞赛中错误率和模型层数的演变：随着模型变深，错误率持续下降\protect\footnotemark}
\end{figure}
\footnotetext{Slides 第21页。}

ImageNet 大规模视觉识别挑战赛（ILSVRC）是推动 CNN 架构发展的重要驱动力。一个清晰的趋势是：随着模型层数的增加和架构设计的改进，错误率持续下降，直到超越人类水平（$\sim$5\%）。

\subsection{AlexNet vs VGGNet}

\begin{figure}[H]
\centering
\includegraphics[width=0.95\textwidth]{slides-images/slide-022.jpg}
\caption{AlexNet 架构：5 个卷积层 + 3 个全连接层\protect\footnotemark}
\end{figure}
\footnotetext{Slides 第23页。}

AlexNet（2012）使用了大小不一的卷积核（11x11, 5x5, 3x3），总共约 6000 万参数。VGGNet（2014）提出了一个重要的设计原则：

\begin{importantbox}{VGGNet 的核心洞察：小滤波器更好}
VGGNet 完全使用 $3 \times 3$ 卷积核代替大卷积核。关键发现：

\textbf{三个 $3 \times 3$ 卷积 $\approx$ 一个 $7 \times 7$ 卷积}（相同的感受野），但：
\begin{itemize}
    \item \textbf{参数更少}：$3 \times (3^2 C^2) = 27C^2$ vs $7^2 C^2 = 49C^2$，减少 45\%
    \item \textbf{表达能力更强}：三层之间有激活函数，引入了更多非线性
    \item \textbf{更深的网络}：更多层 = 更复杂的特征变换
\end{itemize}
\end{importantbox}

\begin{figure}[H]
\centering
\includegraphics[width=0.95\textwidth]{slides-images/slide-024.jpg}
\caption{VGGNet 架构：全部使用 $3 \times 3$ 卷积，层数 16 或 19 层\protect\footnotemark}
\end{figure}
\footnotetext{Slides 第25页。}

\begin{figure}[H]
\centering
\includegraphics[width=0.95\textwidth]{slides-images/slide-025.jpg}
\caption{$3 \times 3$ 卷积 vs 大卷积核的参数对比\protect\footnotemark}
\end{figure}
\footnotetext{Slides 第26页。}

VGGNet 的架构设计模式清晰而规律：
\begin{itemize}
    \item 所有卷积层：$3 \times 3$，stride 1，padding 1（保持空间尺寸）
    \item 池化：$2 \times 2$ max pooling，stride 2（空间尺寸减半）
    \item 每次池化后通道数翻倍：$64 \rightarrow 128 \rightarrow 256 \rightarrow 512$
\end{itemize}

\subsection{本章小结}

从 AlexNet 到 VGGNet，CNN 架构设计的趋势是：用更小的滤波器和更深的网络取代大滤波器和浅网络。小滤波器参数少、非线性更强，但需要更多层才能达到相同的感受野。VGG 确立了``$3 \times 3$ 卷积 + 每阶段下采样并翻倍通道数''的设计范式。

%% ========== Part 6: ResNet ==========

